\section{问题一：安全运输能力与货箱组批}
\label{sec:q1}

\subsection{问题分析与求解思路}
\subsubsection{单点往返的运输能力限制}
本问每个架次只执行$O01\to S_i\to O01$，不跨服务区组批，也不决定实体飞机和电池的并行日程。需要先回答同一机型面对不同地形与航程能安全携带多少质量，再决定完整货箱如何组合。额定载荷只是一个上限，不能忽略空载返航的能量需求后直接作为安全载荷。

\subsubsection{从安全载荷到离散组批}
先计算15个服务区与3类机型的45组安全载荷，再在各区枚举非空货箱组合，筛去质量、体积或能量不合格的组合。可行组合构成有限批型集，后续选择批型使每箱恰好运输一次。本文先最少架次，再比较能耗，最后比较累计作业时间；另保留不同目标向量的非支配方案，使优先关系可见而不依赖任意权重。

\subsection{安全载荷与组批模型}
\subsubsection{最大安全载荷计算}
对服务区$i$与机型$m$，记携带质量$w$的完整往返能耗为$E_{im}(w)$。安全载荷定义为
\begin{equation}
 w_{im}^{\mathrm{safe}}=\max\{w\in[0,Q_m]:E_{im}(w)\le(1-\alpha_m)E_m\}.
\label{eq:q1_safe_payload}
\end{equation}
在给定模型中，载荷增加使等效航程下降，水平和爬升能耗均不减。先检查额定载荷；额定载荷不满足而空载可行时，对边界方程$E_{im}(w)=(1-\alpha_m)E_m$求根。空载亦不可行时该组合没有可用运输能力，而不是将不成立的根强行赋为零。保存结果以Brent法与独立二分法核对。

基准20\%余量下，45项中39项达到额定质量，6项先受能量限制。图\ref{fig:q1_payload}以颜色表示额定能力保留比例，格内给出公斤数，避免将大机型的绝对载荷与相对折减混淆。S008的B、C型安全载荷分别为28.801、58.904\,kg；C型额定80\,kg不能直接用于该区。
\fig[.78]{q1_payload.pdf}{20\%返航安全余量下的最大安全载荷。格内为kg，颜色为安全载荷与额定载荷之比；描框项受到能量限制。}{fig:q1_payload}

\subsubsection{可行货箱组合与完整配送约束}
对服务区$i$的箱集$\mathcal B_i$，任一非空子集$S\subseteq\mathcal B_i$与机型$m$若满足
\begin{equation}
 \sum_{b\in S}\mu_b\le Q_m,\qquad
 \sum_{b\in S}\nu_b\le V_m,\qquad
 E_{im}\!\left(\sum_{b\in S}\mu_b\right)\le(1-\alpha_m)E_m,
\label{eq:q1_pattern_feasibility}
\end{equation}
则形成可行批型$p=(i,m,S)$，其能耗和作业时间为$E_p,T_p$。记全集$\mathcal P=\bigcup_i\mathcal P_i$，变量$y_p\in\{0,1\}$表示是否采用该批型，完整配送要求
\begin{equation}
 \sum_{p:b\in S_p}y_p=1,\qquad \forall b\in\mathcal B.
\label{eq:q1_cover}
\end{equation}
每箱恰好被一个批型覆盖，同时禁止跨服务区组合。质量安全载荷只辅助理解能力边界；生成批型时仍独立检查体积和完整能耗。

\subsubsection{架次、能耗与累计时间的比较顺序}
选择的目标为
\begin{equation}
 \min_{\mathrm{lex}}\left(N,E,T^{\mathrm{sum}}\right)
 =\min_{\mathrm{lex}}\left(\sum_py_p,\sum_pE_py_p,\sum_pT_py_p\right).
\label{eq:q1_objective}
\end{equation}
实现时先求$N^*$，固定$N=N^*$后求最低能耗，再固定前两项求累计时间。该顺序是本文明确采用的决策偏好；不能因此声称题目只有这一种多目标解释。非支配方案用于说明改变偏好会带来什么代价。

\subsubsection{按服务区分解的等价性}
问题一禁止跨区组批，且不计实体并行与电池接续，故不同服务区之间没有共享时序约束。设各区可行组批集合为$\mathcal X_i$，则
\begin{equation}
\mathcal X=\mathcal X_1\times\cdots\times\mathcal X_{15},\qquad
(N,E,T^{\mathrm{sum}})=\sum_i(N_i,E_i,T_i).
\label{eq:q1_separable}
\end{equation}
若一个全局最少架次方案在某区使用了多于局部最少的架次数，用该区局部最少方案替换即可减少总架次且不影响其他区，产生矛盾。因此$N^*=\sum_iN_i^*$。固定各区最少架次后，同样的替换论证适用于能耗；最后固定前两层，累计时间也可逐区最小化。

这种分解来自约束可分离性，而非启发式近似。问题二允许跨区并使用有限实体，任务之间出现资源耦合，因而不再具有式\eqref{eq:q1_separable}的乘积结构。

\subsubsection{容量下界与候选筛选}
对需求质量$M_i$、体积$V_i^{dem}$，可构造必要下界
\begin{equation}
N_i\ge\max\left\{1,\left\lceil\frac{M_i}{Q^{max}}\right\rceil,
\left\lceil\frac{V_i^{dem}}{V^{max}}\right\rceil\right\}.
\label{eq:q1_capacity_lower}
\end{equation}
质量分母还可替换为该区全部机型安全载荷的最大值。质量界与体积界分别是必要条件，其最大值仍有效；是否达到它则由实际不可拆货箱覆盖决定。S001、S002和S003超过80 kg形成三项两架次下界，其余12区至少一架次，总下界为18。

候选筛选依次检查目的地与箱号、质量与体积、完整任务能源。便宜的组合条件先执行，避免对必定超载的箱组重复查询物理代价。在同一区域同一机型下，增加箱子只会增加质量、体积和给定单调能耗，因此已经超限的组合不能通过继续加箱恢复。筛选后保留原箱号，以便从计数解恢复完整配送清单。


\subsection{精确求解方法}
\subsubsection{原始货箱模型与等价货箱压缩}
箱级整数模型直接对原始箱号生成批型，适合核对完整配送。若同区货箱在质量、体积和本问相关属性上完全一致，箱号置换不改变物理可行性与目标，可用每类箱数压缩对称性。压缩只减少重复代表，并未把不可拆箱改成可连续分割的质量。求得计数方案后再恢复唯一箱号，检查80箱无遗漏和重复。

\subsubsection{剩余需求状态的动态规划}
将某区各等价货箱类别的剩余数量记为$\boldsymbol n=(n_1,\ldots,n_K)$，批型$p$携带计数向量$\boldsymbol c_p$。递推为
\begin{equation}
 F(\boldsymbol n)=\min_{\substack{p:\boldsymbol c_p\le\boldsymbol n}}^{\mathrm{lex}}
 \left\{F(\boldsymbol n-\boldsymbol c_p)+(1,E_p,T_p)\right\},\qquad F(\boldsymbol0)=(0,0,0).
\label{eq:q1_dp}
\end{equation}
每次至少减少一个箱子，状态按剩余总数递增计算，不存在循环。不同服务区互不共享批型且目标可加，总体词典序解由各区最优结果相加得到。需要Pareto集合时，每个状态保留未被架次、能耗、时间三项同时支配的标签，再逐区合并并删除支配标签。

\subsubsection{递推顺序、支配筛选与原箱号恢复}
某区有$K_i$个等价类，需求计数为$d_{ik}$，计数状态数量为
\begin{equation}
|\mathcal N_i|=\prod_{k=1}^{K_i}(d_{ik}+1).
\label{eq:q1_dp_states}
\end{equation}
每个非空批型至少减少一箱，故按箱数递增计算状态时，前驱已经处理。原点为$(0,0,0)$，不可达状态保持空标记，不能被默认填成零成本。直接枚举状态与批型的实现时间为$O(|\mathcal N_i||\mathcal P_i|)$，单目标词典序标签及前驱需要$O(|\mathcal N_i|)$存储。

每次更新同时保存前驱状态和所选批型。最终从完整需求状态回溯，获得机型和计数批型；再从相应等价类的原始箱号池中取出对应数量，恢复唯一箱号。恢复后按原质量、体积和完整往返能耗重新核验，计数压缩不会把货箱变成可分割的连续质量。

多目标递推中只在相同剩余需求状态内比较支配关系。若标签$u$在架次、能耗、时间均不大于$v$，则它们接上相同后续批型后仍保持支配，故$v$可删除；剩余需求不同的状态不能这样互相删除。各区标签集再逐次求和并筛除支配目标。Bouman等在无人机路径研究中使用动态规划压缩重复决策\cite{bouman2018}，本文采用的状态则直接对应本题的不可拆需求计数。

\begin{table}[htbp]\centering\small
\caption{精确组批递推与方案恢复步骤}\label{tab:dp_execution}
\begin{tabularx}{\linewidth}{p{1.5cm}L}\toprule
阶段 & 操作\\\midrule
输入 & 读取等价类需求及质量、体积、能源均合格的计数批型。\\
初始化 & 原点赋零标签，其他状态不可达；设置前驱记录。\\
递推 & 按箱数递增遍历，尝试不超过当前需求的批型，累加真实成本。\\
筛选 & 词典序求解保留最优标签；多目标求解保留同状态非支配标签。\\
恢复 & 回溯机型与批型，从原箱号池中分配完整货箱。\\
核验 & 核对箱号全集、容量、返航能量及与整数模型的目标一致性。\\\bottomrule
\end{tabularx}
\end{table}


\subsubsection{多种精确方法的交叉核验}
基准容量筛选产生742个候选计数批型，完整能量检查后保留732个。已有箱级整数规划、计数压缩整数规划和状态动态规划得到相同目标。本轮另以独立物理计算和计数状态递推重新求解基准与扰动情景，基准再次得到18架次、59.131296\,kWh、546.266929\,min。交叉一致说明实现之间对应，但最少架次的证明还需要下一节的解析下界，而非只靠多个程序输出相同值。

\subsection{运输方案与结果分析}
\subsubsection{18架次方案及最少架次证明}
最大额定载荷为80\,kg。S001需求154\,kg，S002、S003各81\,kg，因此三者各至少两架次，其他12区至少一架次：
\begin{equation}
 N\ge2+2+2+12=18.
\label{eq:q1_sortie_lb}
\end{equation}
找到完整可行的18架次方案后，上下界闭合，故最少架次数为18。在固定18架次的条件下，正式方案总能耗59.131296\,kWh，累计作业时间546.266929\,min，A/B/C架次数为0/9/9。A型未被该目标顺序选中，不代表其不能执行任何任务。

\subsubsection{服务区组批与机型分配}
表\ref{tab:q1_area_plan}给出全部区域安排；附录\ref{app:full_schedules}保留18个架次的箱号和属性。S001为两个C型架次，S002和S003各采用一个B型与一个C型架次，其余12区各一架次，与下界结构对应。每项结果是箱组分配，不是货舱中的三维摆放方案。
\begin{table}[htbp]\centering\small
\caption{问题一各服务区正式运输安排}\label{tab:q1_area_plan}
\begin{tabular}{@{}lclrrr@{}}
\toprule
区域 & 架次 & 机型 & 质量/kg & 能耗/kWh & 累计时间/min\\
\midrule
S001 & 2 & C+C & 154 & 5.913 & 53.106\\
S002 & 2 & B+C & 81 & 8.435 & 64.308\\
S003 & 2 & B+C & 81 & 8.545 & 74.177\\
S004 & 1 & C & 59 & 6.153 & 38.351\\
S005 & 1 & C & 59 & 4.222 & 31.287\\
S006 & 1 & C & 59 & 2.598 & 26.026\\
S007 & 1 & C & 45 & 3.411 & 31.793\\
S008 & 1 & C & 45 & 5.836 & 36.662\\
S009 & 1 & B & 25 & 2.096 & 25.974\\
S010 & 1 & B & 25 & 1.823 & 25.923\\
S011 & 1 & B & 25 & 1.074 & 19.807\\
S012 & 1 & B & 25 & 2.752 & 32.046\\
S013 & 1 & B & 25 & 1.825 & 25.175\\
S014 & 1 & B & 25 & 2.140 & 31.163\\
S015 & 1 & B & 25 & 2.307 & 30.471\\
合计 & 18 & B/C=9/9 & 758 & 59.131296 & 546.266929\\
\bottomrule
\end{tabular}
\end{table}


\subsubsection{基线方案与架次—能耗权衡}
在相同输入与20\%余量下，单机型精确解与质量优先启发式对比如表\ref{tab:q1_baselines}。混合精确方案相对仅C型精确方案节能约20.02\%，相对本次混合机型首次适应递减（FFD）基线节能约21.28\%。前一种比较改变机型可选范围，后一种比较求解策略；不能把两类收益都归因于同一算法技巧。
\begin{table}[htbp]\centering\small
\caption{单机型、混合机型与启发式基线}\label{tab:q1_baselines}
\begin{tabular}{@{}llrrr@{}}
\toprule
机型范围 & 求解方法 & 架次 & 能耗/kWh & 累计时间/min\\
\midrule
仅A型 & 精确 & 52 & 112.180153 & 1488.763007\\
仅B型 & 精确 & 35 & 68.350657 & 927.625473\\
仅C型 & 精确 & 18 & 73.934908 & 563.979235\\
混合机型 & 精确 & 18 & 59.131296 & 546.266929\\
混合机型 & 质量优先FFD & 18 & 75.116697 & 563.979235\\
\bottomrule
\end{tabular}
\end{table}


完整离散非支配目标有两组：18架次方案的能耗与累计时间如前所述；19架次方案的能耗为59.033942 kWh、累计时间为573.999623 min。增加一架次节省0.097354\,kWh，约0.165\%，却增加27.733\,min累计时间，变化来自S008由一个C型架次改为两个B型架次。本文因此保留18架次作为主方案，同时保留19架次方案说明最少架次与最低能耗并不等价。

\subsubsection{混合机型的成本分担}
表\ref{tab:rev_q1_types}按机型分解18架次正式方案。B型与C型各9架次，但分别承担225 kg和533 kg。相同架次数并不表示相同任务质量或能源成本；混合机队按服务区需求和具体航段选择机型，而非平均分配物资。
\begin{table}[htbp]\centering\small
\caption{18架次正式组批方案的机型分工}\label{tab:rev_q1_types}
\begin{tabular}{lrrrr}\toprule
机型 & 架次 & 质量/kg & 能耗/kWh & 累计作业/min\\\midrule
A & 0 & 0 & 0.000000 & 0.000\\
B & 9 & 225 & 19.511736 & 255.512\\
C & 9 & 533 & 39.619560 & 290.755\\
\bottomrule\end{tabular}
\end{table}


仅C型精确方案同样达到18架次，因而可作为控制架次数的对照。混合方案减少14.803612 kWh能耗、17.712306 min累计作业时间，说明优化收益并非只是少飞几次，而来自同架次数下的机型与箱组选择。质量优先FFD也产生18架次，却没有获得相同的次级成本，进一步说明完整物理成本需要进入组合选择。

18与19架次方案之间的差分对应每节省1 kWh增加约284.86 min累计作业时间。这个比例用于解释两份离散方案的选择代价，反映S008由一项C型任务改成两项B型任务后，准备、交接和往返工作量的变化。基准架次优先规则选择18架次，19架次则保留为偏重能源目标的可执行备选。

\subsubsection{参数变化与实际任务调整}
对批型$p$定义$\rho_p=1-E_p/E_m$，其安全余量可行条件是$\alpha\le\rho_p$。因此候选集合仅在有限临界值处变化，连续参数能够通过不可拆组合产生离散决策跃迁。事件枚举据此覆盖整个参数范围，比仅比较若干网格点更能定位任务调整时机。

“最少架次数不变”和“原任务无需改变”应分别判断。例如28\%与32\%场景均为20架次，但机型构成不同，实际箱组仍需更新。基准原方案在23.08835\%以内保持可行，则意味着在这一安全要求范围内可以直接沿用已发出的任务清单。超过首个阈值时先复核S004，随后关注S008，使参数分析对应具体服务区与执行操作。


\subsection{稳健性与敏感性分析}
\label{sec:q1_robust}
\subsubsection{原组批方案的安全余量容忍范围}
固定18架次的箱组和机型，统一提高返航余量。原方案允许的最大余量为各架次返航SOC的最小值：
\begin{equation}
 \alpha_{\mathrm{keep}}=\min_r SOC_r^{\mathrm{return}}=0.2308834976366333.
\label{eq:q1_keep}
\end{equation}
限制任务是S004的59\,kg C型架次。因此在$20\%\le\alpha\le23.08834976366333\%$时，原方案无需改变即可使用。进一步，因提高余量只会缩小可行域，原基准词典序最优解仍在其中，所以它在该区间继续最优。该论证依赖其他输入和目标不变，不推广到重量、地形或效率同时变化。

\subsubsection{安全余量对载荷能力与最少架次的影响}
\textbf{最大安全载荷的参数响应。}由式\eqref{eq:q1_safe_payload}，$\alpha$提高只会缩小能量预算，因此各机型的$w_{im}^{\mathrm{safe}}(\alpha)$单调不增。在额定质量处仍有能量富余时，曲线先保持平台；能量限制生效后才下降。为直接回应运输能力的敏感性要求，表\ref{tab:q1_payload_response}列出S001、S004和S008的代表值；全部45个组合的44个余量水平保留在配套数据中。基准20\%结果逐项与原始45项结果核对，最大差小于$10^{-7}$\,kg。
\begin{table}[htbp]\centering\small
\caption{安全余量变化下代表服务区的最大安全载荷/kg}\label{tab:q1_payload_response}
\begin{tabular}{lrrrr}\toprule 服务区 & 安全余量/\% & A型 & B型 & C型\\\midrule
S001 & 20.0000 & 25.000 & 30.000 & 80.000\\
S001 & 30.0000 & 25.000 & 30.000 & 80.000\\
S001 & 35.2678 & 25.000 & 30.000 & 80.000\\
S004 & 20.0000 & 25.000 & 30.000 & 63.697\\
S004 & 30.0000 & 20.169 & 23.442 & 44.747\\
S004 & 35.2678 & 3.591 & 17.455 & 27.362\\
S008 & 20.0000 & 25.000 & 28.801 & 58.904\\
S008 & 30.0000 & 14.213 & 21.001 & 36.814\\
S008 & 35.2678 & --- & 13.998 & 14.000\\
\bottomrule\end{tabular}
\par\smallskip\footnotesize 临界余量计算使用未舍入值；表中35.2678\%仅为显示值。载荷是质量上限，组批还需检查货箱体积。
\end{table}

\fig[.92]{q1_payload_response_S008.pdf}{S008三类机型的安全载荷随返航余量变化。水平虚线为不可拆14\,kg饮用水箱，末端临界点由连续物理函数求根，未用1\%网格插值替代。}{fig:q1_payload_response}
图\ref{fig:q1_payload_response}中，三类机型在临界余量附近的可行载荷与14\,kg单箱相比较，解释了为什么最终不是增加几个架次就能恢复配送：一旦该箱对所有机型都超出安全能力，完整覆盖即不可行。安全载荷下降并不总立刻改变整数架次数，因此需要同时观察以下组批响应。

允许重新组批时，已有完整参数剖面包含569个候选事件，其中8个改变全局最少架次数。图\ref{fig:q1_reserve}给出阶梯变化。阈值本身仍属于左侧较少架次的可行区间，严格越过后才跳升；当$\alpha>35.267806887191433\%$时，S008的一只14\,kg饮用水箱不能被任何机型安全送达，增加架次也无法补救。
\fig[.91]{q1_reserve.pdf}{返航安全余量与最少架次数。阈值处保留较少架次，严格右侧跳升；阴影为不可行情景，不将其编码为某个有限架次数。}{fig:q1_reserve}

本轮重新求解13个安全余量情景，并保留全部箱组。结果在23.5\%时变为19架次，28\%时为20架次，34\%时为22架次，35\%时为25架次；35.3\%时精确递推在S008无可行覆盖。即使架次数不变，机型构成与次级目标也可能改变，例如28\%和32\%均为20架次，但分别采用C/B为10/10和12/8。平坦的架次阶梯不表示全部决策细节不变。

\subsubsection{能耗估计偏差与组批结构变化}
设模型估计的各任务能耗统一增加$\varepsilon$，而电池容量、路线、载荷和时间参数不变，即$E_r'=(1+\varepsilon)E_r$。这是一种能耗估计压力模型，不直接等同于风场或电池老化。固定原方案的能量容忍边界为
\begin{equation}
 \varepsilon_{\max}=\min_r\frac{(1-\alpha)E_{m_r}}{E_r}-1\approx4.015451\%.
\label{eq:q1_energy_limit}
\end{equation}
同时，该能量约束等价于将安全余量改为$\alpha_{\mathrm{eff}}=1-(1-\alpha)/(1+\varepsilon)$。本轮不仅查询已有剖面，还重新执行精确组批，验证0、1、2、4\%偏差仍需18架次，5\%需19架次，10\%需20架次；新的任务构成与能耗见表\ref{tab:q1_sensitivity}。
\begin{table}[htbp]\centering\small
\caption{统一能耗正偏差下重新求解的最少架次方案}\label{tab:q1_sensitivity}
\begin{tabular}{@{}lrrr@{}}
\toprule
能耗偏差 & 架次 & 扰动后能耗/kWh & 累计时间/min\\
\midrule
0\% & 18 & 59.131296 & 546.267\\
1\% & 18 & 59.722609 & 546.267\\
2\% & 18 & 60.313922 & 546.267\\
4\% & 18 & 61.496548 & 546.267\\
5\% & 19 & 64.127057 & 576.094\\
10\% & 20 & 69.619150 & 605.592\\
\bottomrule
\end{tabular}
\end{table}


另将全部机型爬升效率乘以0.8、0.9、1.0、1.1、1.2，并保持其他参数不变。五个情景均为18架次、B/C各9架次，目标能耗在58.705340—59.770231\,kWh变化。这里效率扰动只作用于爬升能耗，因而比统一增加全部任务能耗的影响小。上述范围内的稳定不应被表述为对任意气象或装备误差不敏感。
